\chapter{Estructura entera de los lectores de constantes}
\label{chap:canales-enteros-constantes}
\index{constantes matemáticas!canales enteros}
\index{raíz digital!residuo y acarreo}
\index{semicorona}

Las generaciones anteriores distinguen la función de cada constante. Este
bloque demuestra la estructura entera de sus lectores. Las cifras
\(729,271,315\) y \(161\) se construyen como \emph{canales enteros relativos
a reglas tipadas declaradas}; no entran como aproximaciones decimales. Sus
estatutos no son idénticos: \(729\) es la potencia cuadrática del canal
\(27\); \(271\) y \(315\) reciben una realización nonádica construida desde
la partición de fase TPK publicada y una segunda realización relativa a la elevación de
Steiner; y
\(161\) queda, en el alcance presente, como identidad entera sin una
realización conjuntista propia. Después se clasifica el catálogo mediante
predicados y calibres publicados. Sólo al final puede actuar una carta
arquimediana que produzca desarrollos decimales.

Este orden separa dos operaciones distintas. Buscar el entero \(271\)
porque recuerda a los primeros dígitos de \(e\) sería una lectura
retrospectiva. Encontrar una región de TPK y llamarla \(e\) sin construir el
tipo de retorno que representa sería una denominación sin mapa. Por ello se
separan los canales enteros, la clasificación orbital condicionada y el
diccionario que los relaciona. El canal entero aislado no selecciona una
constante; la selección pertenece al lector TPK enriquecido, que incorpora
la órbita, la orientación, la supervivencia y la memoria de frontera sin
recibir como entrada el valor decimal que finalmente evalúa.

\section{Residuo visible, raíz digital y memoria de corona}

Para \(n>0\), la raíz digital nonádica es la sección
\[
 \operatorname{dr}_9(n)
 =\begin{cases}
 9,&n\equiv0\pmod9,\\
 n\bmod9,&n\not\equiv0\pmod9.
 \end{cases}
\]
La reducción módulo nueve conserva la clase visible; la sección por
\(\{1,\ldots,9\}\) conserva además la elección del representante nulo.
Definimos el cociente levantado
\[
 q_9^\sharp(n):=\frac{n-\operatorname{dr}_9(n)}9
 \in\mathbb Z_{\geq0}.
\]
Se obtiene la descomposición exacta, que no es la convención usual de resto
en \(\{0,\ldots,8\}\),
\[
 n=9q_9^\sharp(n)+\operatorname{dr}_9(n).
\]
Ésta separa dos lecturas: \(\operatorname{dr}_9(n)\) registra el valor
visible y \(q_9^\sharp(n)\) registra cuántas coronas completas preceden al
representante visible elegido. Ésta es la forma exacta de la imagen
coloquial según la cual el módulo compacta y el acarreo conserva lo que la
compactación perdería.

En APP las hojas aditiva y multiplicativa producen
\[
 \sum\Sigma=405,
 \qquad
 \sum\Pi=459,
 \qquad
 \Delta:=\sum\Pi-\sum\Sigma=54.
\]
La suma de los cocientes aditivos es
\[
 S:=\sum Q^+=45.
\]
Definimos
\[
 \kappa=\frac{\Delta}{2}=27,
 \qquad
 \Theta=5\kappa=135.
\]
Todos estos números proceden del ortograma y de sus dos evaluaciones. No se
ha usado todavía ninguna constante.

El bloque central
\[
 B_c=\begin{pmatrix}7&2\\2&7\end{pmatrix}
\]
produce cuatro palabras ordenadas:
\[
 72,\qquad27,\qquad77,\qquad22.
\]
Satisfacen
\[
 72+27=77+22=99,
 \qquad
 72-27=45,
 \qquad
 77-22=55=54+1.
\]
El centro no permanece aislado: sus dos lecturas complementarias se
proyectan sobre el borde \(99\), mientras sus diferencias recuperan la
corona aditiva \(45\) y el defecto apuntado \(54+1\).

\section{Teorema de los cuatro canales enteros}

\begin{theorem}[Identidades tipadas de los cuatro canales]
\label{thm:derivacion-canales-enteros}
Relativamente a las reglas de composición
\(\kappa\mapsto\kappa^2\), \(\Theta\mapsto2\Theta+1\),
\((\Theta,S)\mapsto2\Theta+S\) y
\((\Theta,\kappa)\mapsto\Theta+\kappa-1\), se definen los cuatro canales
enteros
\[
\begin{aligned}
 C_\alpha&=\kappa^2=729,\\
 C_e&=2\Theta+1=271,\\
 C_\pi&=2\Theta+S=315,\\
 C_\varphi&=\Theta+\kappa-1=161.
\end{aligned}
\]
Además,
\[
\begin{aligned}
 C_e&=5\Delta+1=10\cdot27+1
      =4\Delta+(\Delta+1),\\
 C_\pi&=5\Delta+S=7S,\\
 C_\varphi&=3\Delta-1
      =4\Delta-(\Delta+1),
\end{aligned}
\]
y se cumplen los cierres
\[
 \boxed{1000=C_\alpha+C_e=729+271}
\]
y
\[
 \boxed{C_\pi+C_e+C_\varphi-3=744=729+15.}
\]
El canal de \(\alpha\) también posee la lectura central
\(729=10\cdot72+9\), y el canal de \(e\), la lectura
\(271=10\cdot27+1\).
\end{theorem}

\begin{proof}
Sustituir \(\Delta=54\), \(S=45\), \(\kappa=27\) y
\(\Theta=135\) da las cuatro fórmulas. Las expresiones alternativas son
identidades enteras. La primera completación se obtiene de
\(27^2+2(5\cdot27)+1=729+270+1=1000\); la segunda, de
\(315+271+161-3=744=729+15\). Las dos divisiones decimales se siguen de
las palabras centrales \((72,27)\) y de sus restos ya fijados por la
completación cúbica.
\end{proof}

La igualdad \(1000-3^6=271\) es universal para ese ambiente cúbico. Por sí
sola no distingue a \(e\): en un ambiente de cardinal \(B\), construido con
un alfabeto de cardinal \(q\) en dimensión \(d\), el complemento formal es
\(B-q^d\). La información adicional de esta sección es que \(271\) admite
también una realización como retorno apuntado y la lectura central \(27|1\).
Su asociación con la órbita denominada \(e\) pertenece al diccionario tipado
posterior, no a la resta aislada.

\section{Semicorona nonádica relativa a la regla de fase TPK}

La regla mínima de movimiento del emisor TPK no trata los nueve instantes de
una ventana como una lista intercambiable. Relativamente al origen de fase
publicado, define
\[
 D_9=\{1,\ldots,9\},
 \qquad
 H_+=\{1,4,7\},
 \qquad
 H_-=\{2,5,8\},
 \qquad
 H_0=\{3,6,9\}.
\]
Los instantes de \(H_+\) activan la lectura aditiva, los de \(H_-\) la
lectura multiplicativa y los de \(H_0\) conservan el canal neutral. Pongamos
\[
 H_{\rm act}=H_+\sqcup H_-,
 \qquad
 O_{\rm ph}=\{1,\ldots,8\},
 \qquad
 \star_{\rm ph}=9,
\]
y sea \(P_2(H_{\rm act})\) el conjunto de pares no ordenados de fases
activas. Definimos
\[
\begin{aligned}
 \mathcal F^{\rm ph}_6&=H_{\rm act}\times P_2(H_{\rm act}),\\
 \mathcal F^{\rm ph}_8&=O_{\rm ph}\times P_2(H_{\rm act}),\\
 \Theta_{\rm ph}&=D_9\times P_2(H_{\rm act}).
\end{aligned}
\]

\begin{theorem}[Semicorona relativa a la fase TPK]
\label{thm:semicorona-fase-tpk}
Relativamente al origen de ventana publicado y al retorno terminal
\(9\to1\), se cumple
\[
 |\mathcal F^{\rm ph}_6|=90,
 \qquad
 |\mathcal F^{\rm ph}_8|=120,
 \qquad
 |\Theta_{\rm ph}|=135.
\]
La frontera orientada, el retorno apuntado, la frontera lateral neutral y el
cierre lateral son, respectivamente,
\[
\begin{aligned}
 \partial_{\rm ph}&=\Theta_{\rm ph}\times\{-1,+1\},\\
 E_{\rm ph}&=\partial_{\rm ph}\sqcup\{\infty\},\\
 L_{\rm ph}&=H_0\times P_2(H_{\rm act}),\\
 P_{\rm ph}&=\partial_{\rm ph}\sqcup L_{\rm ph},
\end{aligned}
\]
y satisfacen
\[
 \boxed{|\partial_{\rm ph}|=270,\quad
 |E_{\rm ph}|=271,\quad |L_{\rm ph}|=45,\quad
 |P_{\rm ph}|=315.}
\]
La construcción es covariante bajo un cambio de origen: trasladar la ventana
por \(C_9\) traslada simultáneamente todos los estratos y no cambia sus
cardinales ni sus incidencias.
\end{theorem}

\begin{proof}
La regla de signo da seis fases activas y tres neutrales, de modo que

\[
 |P_2(H_{\rm act})|=\binom62=15.
\]

Multiplicar por \(6\), \(8\) y \(9\) produce \(90\), \(120\) y \(135\).
La orientación duplica \(135\); el retorno apuntado añade un elemento; y las
tres fases neutrales aportan \(3\cdot15=45\) posiciones laterales. La
covariancia se sigue de que una traslación de \(D_9\) conjuga la regla de fase
y transporta cada producto cartesiano por una biyección.
\end{proof}

Este teorema mejora la identidad cardinal \(270+1\) de manera sustantiva: el
orden \(6\subset8\subset9\) ya no se elige entre \(9!\) órdenes posibles,
sino que procede del reloj de fase TPK y de su retorno publicado. La
canonicidad es relativa al origen de la ventana; cambiarlo no altera el
objeto hasta la acción natural de \(C_9\).

\section{Semicorona relativa a la elevación
hexada--octada y retorno apuntado}

Fijemos la dodecada, el alineamiento y una hexada
\(H\in\mathcal H(D)\) de la construcción de
\cref{sec:w12-w24-elevacion}. Por el
\cref{thm:elevacion-unica-hexada}, existe una única octada \(O_H\) tal que
\(O_H\cap D=H\). Escribamos
\[
 P_2(H)=\{\{a,b\}:a,b\in H,\ a\ne b\},
 \qquad |P_2(H)|=\binom62=15,
\]
y añadamos a \(O_H\) un estado de retorno \(\star\):
\[
 \widehat O_H=O_H\sqcup\{\star\},
 \qquad |\widehat O_H|=9.
\]
Definimos
\[
 \Theta_H=\widehat O_H\times P_2(H).
\]
Dentro de él aparecen los dos estratos combinatorios
\[
 \mathcal F_6=H\times P_2(H),
 \qquad
 \mathcal F_8=O_H\times P_2(H).
\]

\begin{theorem}[Realización de la semicorona]
\label{thm:realizacion-semicorona}
Los objetos anteriores satisfacen
\[
 |\mathcal F_6|=90,
 \qquad
 |\mathcal F_8|=120,
 \qquad
 |\Theta_H|=135.
\]
La frontera orientada
\[
 \partial_H=\Theta_H\times\{-1,+1\}
\]
tiene \(270\) estados. La adjunción de un punto de retorno,
\[
 E_H=\partial_H\sqcup\{\infty\},
\]
tiene \(271\) estados. La frontera lateral del estrato de seis puntos es
\[
 L_H=\Theta_H\setminus\mathcal F_6,
 \qquad |L_H|=45,
\]
y el cierre orientado lateral
\[
 P_H=\partial_H\sqcup L_H
\]
tiene \(315\) estados.
\end{theorem}

\begin{proof}
Los cardinales son
\[
 6\binom62=90,
 \qquad
 8\binom62=120,
 \qquad
 9\binom62=135.
\]
La orientación duplica \(135\), la adjunción de un punto añade una unidad y
la diferencia entre \(\Theta_H\) y el estrato de seis puntos contiene
\((9-6)15=45\) elementos. De aquí resultan \(270,271\) y
\(270+45=315\).
\end{proof}

Aquí los dos enteros poseen tipos conjuntistas diferentes. \(271\) es el
cardinal de una frontera orientada a la que se ha adjuntado un punto; no es,
por ese solo hecho, la órbita de \(e\). \(315\) es el cardinal de esa
frontera junto con las \(45\) posiciones laterales que no pertenecen al
estrato de seis puntos. La inclusión \(H\subset O_H\) usada aquí sí es la
elevación de Steiner ya demostrada, relativamente a la dodecada y al
alineamiento fijados. La semicorona añade después el punto \(\star\), la
orientación y la frontera lateral. Esta construcción no proporciona todavía
el orden nonádico inducido por TPK ni un entrelazador entre las caras del cubo
y las incidencias de Steiner; ésos son mapas adicionales.

Las dos semicoronas comparten el mismo perfil cardinal, pero no deben
identificarse por ello. Una isomorfía no marcada entre las cadenas
\[
 H_{\rm act}\subset O_{\rm ph}\subset D_9
 \quad\text{y}\quad
 H\subset O_H\subset\widehat O_H
\]
que preserve el punto de retorno puede permutar libremente los seis elementos
del primer estrato y los dos elementos añadidos en el segundo. Existen
\[
 \boxed{6!\,2!=1440}
\]
calibres de cadena. El entrelazador fase--Steiner debe seleccionar uno de
ellos mediante incidencias adicionales; la igualdad de cardinales no lo
construye.

\section{Acción diedral y clasificación bajo predicados declarados}

Para una emisión
\(U=(u_1,\ldots,u_6)\in\mathcal U_6\), definimos la palabra visible
\[
 \Psi(U)=((-u_1)\bmod3,\ldots,(-u_6)\bmod3)
 \in\mathbb F_3^6.
\]
Sean las permutaciones de coordenadas definidas explícitamente por
\[
 r\cdot(u_1,u_2,u_3,u_4,u_5,u_6)
 =(u_3,u_1,u_2,u_5,u_6,u_4),
\]
\[
 s\cdot(u_1,u_2,u_3,u_4,u_5,u_6)
 =(u_4,u_5,u_6,u_1,u_2,u_3).
\]
Generan
un grupo diedral \(D_3\simeq S_3\). El catálogo verifica
\[
 r\mathcal U_6=\mathcal U_6,
 \qquad
 s\mathcal U_6=\mathcal U_6,
 \qquad
 \Psi(gU)=g\Psi(U),
\]
y conserva las multiplicidades microscópicas.

La acción divide las \(243\) palabras visibles en \(38\) órbitas de tamaño
seis y cinco órbitas de tamaño tres. Sobre cada palabra \(w\), escribamos
\[
 n(w)=|\Psi^{-1}(w)|,
 \qquad
 M(w)=\sum_{\Psi(U)=w}|F^{-1}(U)|.
\]

Fijamos ahora tres predicados publicados. El predicado \(P_\pi\) exige
\(n(w)=5\), \(M(w)=1008\), multiplicidades
\((432,144,144,144,144)\) y censo trítico \((2,3,1)\). El predicado
\(P_e\) exige, sobre la órbita completa, emisiones unipuntuales de masa
microscópica \(144\), unión de valores 
\(\bigcup_{w\in\mathcal O}\operatorname{diag}_c(w)=\{0,3,6\}\) y censo
\((2,3,1)\). El predicado \(P_\varphi\) conserva las
tres primeras condiciones de \(P_e\) y exige censo equilibrado \((2,2,2)\).
Los calibres de orientación son \(\operatorname{diag}_c=6\) y la clase
cardinal \(ES\).

\begin{theorem}[Clasificación orbital finita condicionada]
\label{thm:selectores-orbitales-constantes}
En el catálogo reconstruido, los predicados anteriores seleccionan las
siguientes órbitas únicas sin consultar dígitos de ninguna constante.
\begin{enumerate}
  \item Existe una única órbita de seis palabras para la cual, en cada
  palabra, \(n(w)=5\), \(M(w)=1008\) y las cinco multiplicidades son
  \(432,144,144,144,144\). Todas sus palabras tienen censo trítico
  \((n_0,n_1,n_2)=(2,3,1)\). Es la órbita diedral del tipo regional de
  \(\pi\).
  \item Entre las órbitas radicales de emisiones unipuntuales y masa
  microscópica \(144\) cuya unión orbital de valores
  \(\operatorname{diag}_c\) es \(\{0,3,6\}\), existe una única
  órbita con censo \((2,3,1)\). Es la órbita de retorno de \(e\).
  \item En el mismo sector existe una única órbita equilibrada con censo
  \((2,2,2)\). Es la órbita proporcional de \(\varphi\).
\end{enumerate}
La orientación declarada por
\(\operatorname{diag}_c=6\) y la clase cardinal \(ES\) selecciona,
respectivamente,
\[
 \boxed{
 w_\pi=010211,
 \qquad
 w_e=201101,
 \qquad
 w_\varphi=121200.
 }
\]
\end{theorem}

\begin{proof}
La reconstrucción exhaustiva de las \(468\) emisiones comprueba primero el
cierre de la acción y agrupa las \(243\) palabras en órbitas. Después aplica
los predicados indicados a todas ellas. El primer predicado deja una sola
órbita; los dos predicados de fibra radical unipuntual dejan, respectivamente, una
órbita de censo \((2,3,1)\) y una de censo \((2,2,2)\). Finalmente se
aplica el calibre de orientación a cada órbita y queda una única palabra.
Ninguno de los predicados del selector contiene los desarrollos decimales de
\(\pi,e\) ni \(\varphi\).
\end{proof}

\subsection{Primer codificador intrínseco y su límite}

En el estrato de palabras visibles que contienen exactamente dos ceros existe
un codificador que sí procede del propio dato TPK:
\[
 c_0(w)=\{j\in\{1,\ldots,6\}:w_j=0\}
 \in P_2(\{1,\ldots,6\}).
\]
Es \(D_3\)-equivariante porque una permutación de coordenadas transporta el
soporte de ceros. En las órbitas clasificadas por el
\cref{thm:selectores-orbitales-constantes}, sus imágenes son
\[
\begin{aligned}
 c_0(\mathcal O_\pi)
 &=\{\{1,2\},\{1,3\},\{2,3\},\{4,5\},\{4,6\},\{5,6\}\},\\
 c_0(\mathcal O_e)
 &=\{\{1,6\},\{2,5\},\{3,4\}\},
\end{aligned}
\]
donde en la segunda línea cada par posee dos preimágenes dentro de la órbita.
La órbita de \(\varphi\) comparte el primer tipo de soporte con \(\pi\), pero
se separa de ella por su fibra unipuntual, su multiplicidad y su censo TRIT\@.

El mapa \(c_0\) es el primer peldaño del codificador registro--pares; no es la
familia completa \(c_k\). En bloques profundos aparecen (0,1,3) o (4)
ceros, de modo que la misma fórmula no toma valores en \(P_2\) a todas las
profundidades. Además, \(\{1,\ldots,6\}\) es aquí la hexada de coordenadas de
la palabra, no una hexada de Witt ya identificada. Prolongar \(c_0\) exige
una regla dependiente de fase y el entrelazador de incidencia correspondiente.

La microfibra canónica de \(\pi\) sigue siendo exactamente
\[
 \mathscr R_\pi^{\rm micro}=\Psi^{-1}(010211),
\]
formada por cinco regiones \(U_6\) distintas. Las otras cinco palabras de la
órbita diedral son imágenes de calibre de esa forma orbital; no constituyen
cinco nuevas microfibras canónicas de \(\pi\). Análogamente, los nombres
\(e\) y \(\varphi\) se aplican después de fijar los predicados y el calibre
publicados. La enumeración es exhaustiva \emph{relativamente} a esas
premisas. Predicados y calibre son componentes declaradas del lector TPK
enriquecido: no se reconstruyen a partir de las tres estadísticas gruesas
que aparecen a continuación, sino de la genealogía completa de la emisión,
su orientación, su frontera y su supervivencia.

La suma de los seis enteros de una emisión, su residuo módulo nueve y el
número de entradas pares no bastan para este teorema. La emisión orientada
de \(e\) y tres de las cinco regiones de \(\pi\) comparten exactamente
\[
 \sum_j u_j=3316,
 \qquad
 \sum_j u_j\equiv4\pmod9,
 \qquad
 \#\{j:u_j\text{ par}\}=4.
\]
La distinción exige la estructura de la fibra, la multiplicidad, el sector
radical y la órbita completa. Esta ablación demuestra por qué la paridad y
la raíz digital son ingredientes del lector, pero no sustituyen el
transporte TPK\@.

\section{Correspondencia tipada \texorpdfstring{\(271\to e\) y
\(315\to\pi\)}{271 a e y 315 a pi}}

Los
\cref{thm:semicorona-fase-tpk,thm:realizacion-semicorona,thm:selectores-orbitales-constantes}
son independientes: el primero construye
la semicorona desde el reloj de fase, otro proporciona su realización de
Steiner relativa y el tercero clasifica órbitas bajo predicados fijados. Para
componer el primer y el tercer resultado declaramos el
diccionario de tipos que debe conservar el lector:
\[
\begin{array}{c|c}
\text{tipo interno}&\text{tipo orbital}\\
\hline
\text{retorno apuntado bidireccional}&
\text{fibra radical unipuntual con el censo de la rama regional}\\
\text{cierre de frontera más lateral}&
\text{órbita de cierre con cinco regiones}.
\end{array}
\]

\begin{corollary}[Asociaciones tipadas y orbitales relativas]
\label{cor:lectores-enteros-relativos}
Relativamente al diccionario, a los predicados \(P_e,P_\pi\) y a los
calibres de orientación anteriores, la clasificación orbital produce las
asociaciones condicionadas
\[
 \boxed{|E_{\rm ph}|=271\rightsquigarrow \mathcal O_e
        \longrightarrow w_e,}
 \qquad
 \boxed{|P_{\rm ph}|=315\rightsquigarrow \mathcal O_\pi
        \longrightarrow w_\pi.}
\]
En el segundo caso, el representante orientado \(w_\pi=010211\) determina la
microfibra canónica de cinco regiones \(\Psi^{-1}(010211)\). Las asociaciones
se obtienen antes de la evaluación arquimediana y no usan los valores
decimales de las constantes.
\end{corollary}

\begin{proof}
Cada fila del diccionario posee, por el teorema orbital, una única órbita que
satisface el predicado declarado. La clasificación se efectúa primero al
nivel de la órbita diedral completa; el calibre orientado elige después el
representante indicado. La unicidad es, por tanto, relativa al diccionario,
a los predicados y a los calibres fijados.
\end{proof}

El canal \(161=\Theta+\kappa-1=3\Delta-1\) permanece como identidad entera.
Aunque \(P_\varphi\) clasifica una órbita equilibrada y el calibre publicado
selecciona \(w_\varphi=121200\), esta sección no construye un objeto
conjuntista de cardinal \(161\) cuyo tipo induzca esa órbita. Por ello no se
afirma aquí un lector \(161\to\varphi\).

El corolario depende del diccionario de tipos porque los cardinales, tomados
aisladamente, no definen una función hacia el catálogo. En el lector completo,
ese diccionario se realiza mediante las nueve relaciones de monodromía del
\cref{chap:extensiones-dinamica-cociente}, junto con la orientación y la
supervivencia regional. La consecuencia importante es positiva: una vez
fijado el lector estructural, la evaluación no consulta ninguna aproximación
decimal de la constante.

\subsection{Dependencia conjunta de los canales en el selector de producto}

La factorización ejecutable del catálogo y la parametrización de su firma
aditiva dan
\[
 \mathcal U_6\simeq\mathcal S_+\times\mathcal S_\times,
 \qquad
 \mathcal S_+\simeq D_9\times\{\pm1\}.
\]
Por otra parte,
\[
 \partial_{\rm ph}
 \simeq(D_9\times\{\pm1\})\times P_2(H_{\rm act}).
\]
Así, una tentativa funcional de descenso que conserve de manera idéntica la
fase y la orientación queda reducida al mapa
\[
 q_\times:\mathcal S_\times\longrightarrow P_2(H_{\rm act}).
\]
Éste es el campo ausente del catálogo: las filas publicadas terminan en la
firma producto \(E_{\rm sig}\) y no contienen una clase
\(P_2(H_{\rm act})\).

\begin{proposition}[Obstrucción al selector producto aislado]
\label{prop:nogo-selector-producto-d3}
No existe una función total \(D_3\)-equivariante
\[
 q_\times:\mathcal S_\times\longrightarrow P_2(H_{\rm act})
\]
si \(D_3\) actúa sobre las (26) firmas producto mediante las permutaciones
de coordenadas \(r,s\) publicadas y sobre \(H_{\rm act}\) mediante
\[
 a(x)=x+3\pmod9,
 \qquad
 j(x)=-x\pmod9.
\]
\end{proposition}

\begin{proof}
La acción sobre \(\mathcal S_\times\) tiene órbitas de tamaños
\[
 6+6+6+3+3+1+1.
\]
Sus dos puntos globalmente fijos son las firmas
\(000000\) y \(555555\). La acción de \(D_3\) sobre los seis elementos de
\(H_{\rm act}\) es regular, y su acción inducida sobre los quince pares tiene
órbitas
\[
 6+3+3+3,
\]
sin punto globalmente fijo. Una función equivariante debe enviar un punto
fijo de la fuente a un punto fijo del blanco, lo cual es imposible.
\end{proof}

La obstrucción identifica con precisión la información que pierde la firma
producto aislada: los sectores fijos sólo pueden resolverse conservando fase,
orientación y registro enriquecido. El codificador \(c_0\) anterior actúa en
el estrato de dos ceros, pero no puede sustituir la familia \(c_k\) por una
sola columna de \(E_{\rm sig}\).

\section{Paridad visible, paridad de acarreo y frontera física}

La separación par--impar puede formularse sin depender de una etiqueta de
partícula. Para cada una de las hojas
\(\bullet\in\{+,\times\}\), pongamos
\[
 n_+(i,j)=i+j,
 \qquad
 n_\times(i,j)=ij,
\]
y definamos los dos bits
\[
 \epsilon_R^\bullet(i,j)
 =\operatorname{dr}_9(n_\bullet(i,j))\bmod2,
 \qquad
 \epsilon_Q^\bullet(i,j)
 =q_9^\sharp(n_\bullet(i,j))\bmod2.
\]

\begin{theorem}[Factorización paridad--acarreo]
\label{thm:factorizacion-paridad-acarreo}
Para las \(81\) celdas de cada hoja APP,
\[
 \boxed{
 n_\bullet(i,j)\bmod2
 =\epsilon_R^\bullet(i,j)+\epsilon_Q^\bullet(i,j)\pmod2.}
\]
El censo exacto por pares
\((\epsilon_R,\epsilon_Q)=(00,01,10,11)\) es
\[
 \begin{array}{c|rrrr}
  &00&01&10&11\\
 \hline
 +&16&20&20&25\\
 \times&25&5&20&31
 \end{array}
\]
y, por tanto, las hojas aditiva y multiplicativa poseen distribuciones de
paridad--acarreo distintas.
\end{theorem}

\begin{proof}
Reducir módulo dos
\(n=9q_9^\sharp(n)+\operatorname{dr}_9(n)\) da la identidad, pues
\(9\equiv1\pmod2\). En la hoja aditiva, las multiplicidades de \(i+j\) son
\(1,2,\ldots,9,8,\ldots,1\); separar representante y cociente produce la
primera fila. El mismo cálculo sobre los \(81\) productos \(ij\) produce la
segunda. La enumeración de las \(162\) celdas verifica la identidad antes de
acumular los conteos.\qedhere
\end{proof}

Sobre \(\ell^2(D_9^2)\), los operadores diagonales
\[
 \Gamma_R^\bullet f(i,j)
 =(-1)^{\epsilon_R^\bullet(i,j)}f(i,j),
 \qquad
 \Gamma_Q^\bullet f(i,j)
 =(-1)^{\epsilon_Q^\bullet(i,j)}f(i,j)
\]
son involuciones conmutativas. La paridad sin reducir se factoriza como
\[
 \boxed{\Gamma_{\rm raw}^\bullet
 =\Gamma_R^\bullet\Gamma_Q^\bullet.}
\]
Ésta es la formulación operatorial precisa de la idea de que la raíz digital
lee lo visible mientras el cociente conserva el arrastre de las coronas. La
diferencia entre las dos filas de la tabla es también una forma exacta de la
incompatibilidad entre sumar y acumular producto.

Cada involución descompone el espacio en sectores pares e impares y proporciona
una graduación \(\mathbb Z/2\mathbb Z\) interna. Sin embargo, una estadística
de Bose--Einstein o Fermi--Dirac es una afirmación sobre el intercambio de
varias realizaciones indistinguibles, no sobre la paridad de una única celda.
Su derivación requiere un funtor de intercambio que represente los grupos
simétricos o de trenzas y que entrelace la dinámica. La paridad HMT construye
el dominio graduado sobre el que ese funtor deberá actuar; no se identifica
aquí con su representación física.
